%!TEX root = main.tex

\section{Broadcast Algorithms for MANETs}
\label{sec:background}

In this paper we focus on broadcasting algorithms for mobile ad-hoc networks.
To ease the discussion about the algorithms that are evaluated hereafter, we rely on the taxonomy proposed  by Ruiz et al.~\cite{Ruiz2015}.
In such survey, the authors present many broadcasting algorithms and propose a taxonomy of features to classify them.
In addition, the survey summarizes the terminology used in the literature.
We use the following subset of those terms in the remainder of this paper:

An ad-hoc network can be represented as a graph $G = (V, E)$, where $V$ is the set of
nodes composing it, and $E$ is the set of links connecting nodes in the communication range.

\newtheorem{mydef}{Definition}

\begin{mydef}[Dissemination/Broadcast]
	Given a \textit{source node} $v$ starting a broadcast process, the broadcast problem consists in reaching any $u \in V$ with the message $m$ by following the set of links $e \in E$.
\end{mydef}

\begin{mydef}[Broadcast Session]
	In an ad-hoc network, nodes $V^\prime \subseteq V$ may aim at disseminating the messages $m_0, m_1, \dots, m_k$.
	The dissemination of a message $m_i$ is a \textit{broadcast session}; and it can be uniquely identified by the message.
\end{mydef}

\begin{mydef}[Underlying Topology]
	Dissemination algorithms often maintain a subset of nodes -- the underlying topology -- used to drive the broadcasting process. Those nodes, which play a special role during a broadcast session, are determined using strategies specifics to each dissemination technique.
\end{mydef}
